\chapter{ANÁLISIS REGULATORIO}

Es importante contextualizar el caso Amazon dentro de la línea temporal regulatoria. El desarrollo y abandono del sistema se produjo entre 2014 y 2017, es decir, antes de la entrada en vigor del Reglamento General de Protección de Datos (RGPD), aplicable desde el 25 de mayo de 2018, y muy anterior al Reglamento Europeo de Inteligencia Artificial (AI Act), vigente desde el 1 de agosto de 2024 \cite{rgpd2016,aiact2024}. En Estados Unidos, jurisdicción en la que se desarrolló el sistema, el marco aplicable en aquel momento era el Título VII de la Civil Rights Act de 1964, sin que existieran todavía directrices específicas sobre inteligencia artificial en el empleo.

Si el RGPD hubiera estado plenamente vigente y aplicable al tratamiento, este habría exigido a Amazon identificar una base jurídica clara para el tratamiento automatizado de los currículums, aplicar el principio de minimización de datos y —dado que se trataba de un tratamiento a gran escala con potencial impacto significativo sobre los derechos de las personas (el acceso al empleo)— realizar una Evaluación de Impacto en Protección de Datos (EIPD) antes del despliegue del sistema. Asimismo, el artículo 22 del RGPD, que reconoce el derecho a no ser objeto de una decisión basada únicamente en el tratamiento automatizado que produzca efectos jurídicos o le afecte significativamente, habría sido relevante si el sistema hubiera llegado a tomar decisiones de contratación sin intervención humana real.

En España, la LOPDGDD habría exigido proporcionalidad, información previa a los candidatos y respeto a sus derechos digitales en el contexto de un proceso de selección automatizado \cite{lopdgdd2018}.

El AI Act resulta especialmente pertinente para este análisis: clasifica expresamente los sistemas de inteligencia artificial destinados a la selección o evaluación de candidatos a un empleo como sistemas de alto riesgo. De haber estado vigente esta norma en el momento de los hechos, Amazon habría estado obligada a implementar un sistema de gestión de riesgos, garantizar la calidad y representatividad de los datos de entrenamiento, elaborar documentación técnica detallada, permitir la trazabilidad y auditoría del sistema, y mantener una supervisión humana efectiva y no meramente simbólica; obligaciones que, de haberse cumplido desde el diseño, probablemente habrían forzado la detección del sesgo antes de que el sistema operara durante un año completo sin control \cite{aiact2024}.

En cuanto al marco estadounidense, no fue hasta 2022 y 2023 cuando la EEOC emitió guías técnicas específicas (entre ellas, la relativa a la evaluación del impacto adverso en el uso de software, algoritmos e inteligencia artificial en los procedimientos de selección de empleo bajo el Título VII de la Civil Rights Act de 1964) aclarando que las herramientas algorítmicas de selección de personal están sujetas a las mismas leyes antidiscriminación que cualquier proceso de contratación humano, y que el argumento de que ``el algoritmo lo decidió'' no constituye una defensa válida frente a una denuncia de discriminación \cite{eeoc2023}. Estas guías surgen, en gran medida, como respuesta a casos como el de Amazon.

\section{¿La regulación actual habría evitado este problema?}
Probablemente no lo habría evitado por completo, dado que el sesgo se originó en los datos históricos de contratación y no en una intención discriminatoria explícita ni en una infracción normativa deliberada; ninguna norma puede, por sí sola, ``limpiar'' un histórico social desigual. Sin embargo, la regulación actual sí habría anticipado y mitigado sustancialmente el riesgo: el AI Act habría exigido una auditoría de sesgos y una evaluación de impacto obligatorias antes del despliegue, el RGPD habría reforzado la transparencia y el derecho de los candidatos a comprender la lógica del tratamiento, y las guías de la EEOC habrían dejado claro que Amazon no podría exonerarse de responsabilidad alegando el carácter ``neutral'' del algoritmo. En conjunto, el marco regulatorio actual reduce significativamente la probabilidad de que un caso así permanezca oculto durante un año sin ser corregido, aunque no elimina por completo el riesgo estructural de que los datos históricos reproduzcan desigualdades sociales previas.
